\documentclass[]{exam}
\input{global}
\input{problem-sets}


\setrandomizerseed{1}
\bracketedpoints


\firstpageheader{{Name:\enspace\makebox[6cm]{\hrulefill}}\\Physics}{}{Quiz on Unit 12: Electromagnetic Induction\\}
\runningheader{Physics}{}{Quiz on Unit 12: Electromagnetic Induction}

\begin{document}

\begin{questions}

\question
The lines shown below are called 

% \fbox{
\begin{minipage}{7cm}
\centering
\begin{randomizechoices}
    \correctchoice magnetic field lines
    \choice electric field lines
    \choice gravitational field lines
    \choice tensor field lines
\end{randomizechoices}
\end{minipage}%}%
% \fbox{
\begin{minipage}{6cm}
\centering
\begin{tikzpicture}
    \node at (0,0) {\includegraphics[width=5cm]{documents/figures/magnetic-field-lines-1.pdf}};
\end{tikzpicture}
\end{minipage}%}



\question
What are the labels of the poles of a bar magnet?

\begin{randomizechoices}
\choice North only
\choice North, south, east, and west
\choice East and west
\CorrectChoice North and south
\end{randomizechoices}

\question
What do you get when you cut a bar magnet in half?

\begin{randomizechoices}
\choice One magnet with two south poles, and one magnet with two north poles.
\choice Two objects that are not magnetic anymore.
\choice It’s impossible to cut a magnet in half.
\CorrectChoice Two magnets, each with a north and south pole.
\end{randomizechoices}

\question
True or False? Earth generates a surrounding magnetic field.

\begin{randomizechoices}[norandomize]
    \correctchoice True
    \choice False
\end{randomizechoices}

\question
The south geographic pole of Earth (in Antarctica) is near what type of \textit{magnetic} pole?

\begin{randomizechoices}
    \correctchoice north
    \choice south
    \choice east
    \choice west
\end{randomizechoices}


\question
What type of \textit{magnetic} pole is close to Earth's north \textit{geographic} pole (near Alaska)? 

\begin{randomizechoices}
\choice north
\CorrectChoice south
\choice east
\choice west
\end{randomizechoices}

% \question
% The \textit{magnetic} pole of Earth that is close to the \textit{geographic} South Pole is a\\ magnetic \fillin Pole.

% \begin{randomizechoices}
% \CorrectChoice North
% \choice South
% \choice East
% \choice West
% \end{randomizechoices}

% \question
% The \textit{magnetic} North Pole of Earth is close to the \textit{geographic} \fillin Pole.

% \begin{randomizechoices}
% \choice North
% \CorrectChoice South
% \choice East
% \choice West
% \end{randomizechoices}

\clearpage

\question
Which of the following diagrams correctly shows the Earth's magnetic field lines?

\begin{EnvUplevel}
\centering
\begin{minipage}{4cm}
\centering
\ifprintanswers
    \color{red}
\fi
\begin{tikzpicture}
    \node at (0,0) {\rotatebox{-90}{\fcolorbox{black}{white}{\includegraphics[width=3.5cm]{documents/figures/magnetic-field-lines-1.pdf}}}};
    \node at (0,0) {\twemoji[width=1cm]{globe showing Americas}};
    \node[above] at (current bounding box.north) {Diagram A};
\end{tikzpicture}

\end{minipage}%
\begin{minipage}{4cm}
\centering
\begin{tikzpicture}
    \node at (0,0) {\rotatebox{-180}{\fcolorbox{black}{white}{\includegraphics[width=3.5cm]{documents/figures/magnetic-field-lines-1.pdf}}}};
    \node at (0,0) {\twemoji[width=1cm]{globe showing Americas}};
    \node[above] at (current bounding box.north) {Diagram B};
\end{tikzpicture}
\end{minipage}%
\begin{minipage}{4cm}
\centering
\begin{tikzpicture}
    \node at (0,0) {\rotatebox{-270}{\fcolorbox{black}{white}{\includegraphics[width=3.5cm]{documents/figures/magnetic-field-lines-1.pdf}}}};
    \node at (0,0) {\twemoji[width=1cm]{globe showing Americas}};
    \node[above] at (current bounding box.north) {Diagram C};
\end{tikzpicture}
\end{minipage}%
\begin{minipage}{4cm}
\centering
\begin{tikzpicture}
    \node at (0,0) {\rotatebox{0}{\fcolorbox{black}{white}{\includegraphics[width=3.5cm]{documents/figures/magnetic-field-lines-1.pdf}}}};
    \node at (0,0) {\twemoji[width=1cm]{globe showing Americas}};
    \node[above] at (current bounding box.north) {Diagram D};
\end{tikzpicture}
\end{minipage}%
\end{EnvUplevel}

\scalebox{0}{
\begin{randomizeoneparchoices}[norandomize]
    \correctchoice Diagram A
    \choice Diagram B
    \choice Diagram C
    \choice Diagram D
\end{randomizeoneparchoices}
}


% \question
% The \textit{magnetic} South Pole of Earth is close to the \textit{geographic} \fillin Pole.

% \begin{randomizechoices}
% \CorrectChoice North
% \choice South
% \choice East
% \choice West
% \end{randomizechoices}

\question
One way to visualize magnetic field lines around a magnet is by\ldots 

\begin{randomizechoices}
\choice using certified electromagnetic goggles.
\choice attracting or repelling the magnet with a second magnet.
\CorrectChoice sprinkling iron filings around the magnet.
\choice brushing the magnet with glow-in-the-dark liquid and turning off the lights.
\end{randomizechoices}

\question
Besides magnets, what other thing creates a magnetic field?

\begin{randomizechoices}[norandomize]
\choice Water.
\choice The gravitational force.
\CorrectChoice electric current (i.e., moving charges)
\choice Nothing. Only magnets can create magnetic fields.
\end{randomizechoices}



\question
What is a magnetic field?

\begin{randomizechoices}
\choice a field of magnets
\CorrectChoice directional lines around a magnetic material that indicate the direction and magnitude of the magnetic force
\choice a material or object that produces a magnetism
\choice part of a magnet that exerts the strongest force on other magnets or magnetic material
\end{randomizechoices}


\question
True or false? Monopoles exist in nature.

\begin{randomizechoices}[norandomize]
    \choice True
    \correctchoice False 
\end{randomizechoices}

\question
Why do two north poles of magnets repel each other?

\begin{randomizechoices}
    \correctchoice Their magnetic fields push in opposite directions and do not connect easily.
    \choice North poles contain positive electric charge.
    \choice The magnets lose energy when close together.
    \choice Air pressure forces them apart.
\end{randomizechoices}






\end{questions}
\end{document}